% GENERATED FILE. Edit canonical lessons, then run npm run generate:books.
% canonical-source-hash: fnv1a64:88662f042c1852e8
% canonical-lessons: LA-C358-assentior, LA-C358-dissentio, LA-C358-consentio, LA-C358-consensus, LA-C358-controversia

\chapter{Agreeing and Disagreeing}
\label{ch:la-agreeing-and-disagreeing}
\hlchaptermodality{358}

\begin{chapteropening}
\textbf{By the end of this chapter:} \emph{I can say that I agree or disagree, that we are of one mind, and that a question is in dispute.}

Complete the last lesson of chapter 358: I can say that I agree or disagree, that we are of one mind, and that a question is in dispute.
\end{chapteropening}

\section[assentior, assentīrī]{assentior, assentīrī — to agree (with), to assent}

\label{lesson:LA-C358-assentior}



\begin{quote}
Before the new one: say the Latin for to suppose, then the Latin for an opinion.
\end{quote}

\subsection*{You'll want to know: assentior, assentīrī}
\textbf{assentior, assentīrī} — \textquotedblleft{}to agree (with), to assent\textquotedblright{}.

\begin{quote}
Tibi assentior: nōn est tūtum.
\end{quote}

\begin{quote}
\emph{I agree with you: it is not safe.}
\end{quote}

Deponent: \emph{assentior, assentīrī, assēnsus sum}, with the dative of the person: \emph{tibi assentior}. English \emph{assent}.

\begin{cousinweb}
\emph{assentior} is built from \emph{ad} \textquotedblleft{}towards\textquotedblright{} and \emph{sentiō} \textquotedblleft{}to feel, to think\textquotedblright{} (EtymDB; Lewis \& Short): to bring your feeling over to someone else's. English \emph{assent} comes from it.
\end{cousinweb}

\begin{grammarlens}[title={What you've built}]
One more word for telling, explaining and describing at B1.
\end{grammarlens}

\subsection*{Your turn}
\begin{itemize}
\raggedright
  \item \emph{Say it:} \emph{assentior, assentīrī}
  \item \emph{Say it:} \emph{assentior, assentīrī}, once more
  \item \emph{Say it:} a sentence of your own with \emph{assentior, assentīrī}
  \item \emph{Recall:} say the Latin for to suppose, then the Latin for an opinion, then say \emph{assentior, assentīrī} again
\end{itemize}

\begin{tcolorbox}[breakable,colback=teal!4,colframe=teal!35!black,title={Before you move on}]
What does assentior, assentīrī mean? (\textquotedblleft{}To agree (with), to assent\textquotedblright{}.) Say it once more.
\end{tcolorbox}
\section[dissentiō, dissentīre]{dissentiō, dissentīre — to disagree, to differ in opinion}

\label{lesson:LA-C358-dissentio}



\begin{quote}
Before the new one: say the Latin for an opinion, then the Latin for to agree.
\end{quote}

\subsection*{You'll want to know: dissentiō, dissentīre}
\textbf{dissentiō, dissentīre} — \textquotedblleft{}to disagree, to differ in opinion\textquotedblright{}.

\begin{quote}
Saepe ā frātre meō dissentiō.
\end{quote}

\begin{quote}
\emph{I often disagree with my brother.}
\end{quote}

\emph{dissentiō, dissentīre, dissēnsī, dissēnsum}. The person you disagree with goes in the ablative with \emph{ā / ab}. English \emph{dissent}.

\begin{cousinweb}
\emph{dissentiō} is built from \emph{dis-} \textquotedblleft{}apart\textquotedblright{} and \emph{sentiō} \textquotedblleft{}to feel, to think\textquotedblright{} (EtymDB; Lewis \& Short): to feel apart. English \emph{dissent} comes from it.
\end{cousinweb}

\begin{grammarlens}[title={What you've built}]
One more word for telling, explaining and describing at B1.
\end{grammarlens}

\subsection*{Your turn}
\begin{itemize}
\raggedright
  \item \emph{Say it:} \emph{dissentiō, dissentīre}
  \item \emph{Say it:} \emph{dissentiō, dissentīre}, once more
  \item \emph{Say it:} a sentence of your own with \emph{dissentiō, dissentīre}
  \item \emph{Recall:} say the Latin for an opinion, then the Latin for to agree, then say \emph{dissentiō, dissentīre} again
\end{itemize}

\begin{tcolorbox}[breakable,colback=teal!4,colframe=teal!35!black,title={Before you move on}]
What does dissentiō, dissentīre mean? (\textquotedblleft{}To disagree, to differ in opinion\textquotedblright{}.) Say it once more.
\end{tcolorbox}
\section[cōnsentiō, cōnsentīre]{cōnsentiō, cōnsentīre — to agree together, to be of one mind}

\label{lesson:LA-C358-consentio}



\begin{quote}
Before the new one: say the Latin for to agree, then the Latin for to disagree.
\end{quote}

\subsection*{You'll want to know: cōnsentiō, cōnsentīre}
\textbf{cōnsentiō, cōnsentīre} — \textquotedblleft{}to agree together, to be of one mind\textquotedblright{}.

\begin{quote}
Omnēs cīvēs dē hāc rē cōnsentiunt.
\end{quote}

\begin{quote}
\emph{All the citizens agree about this matter.}
\end{quote}

\emph{cōnsentiō, cōnsentīre, cōnsēnsī, cōnsēnsum}. Where \emph{assentior} is one person agreeing with another, \emph{cōnsentiō} is many minds agreeing together. English \emph{consent}.

\begin{cousinweb}
\emph{cōnsentiō} is built from \emph{con-} \textquotedblleft{}together\textquotedblright{} and \emph{sentiō} \textquotedblleft{}to feel, to think\textquotedblright{} (EtymDB; Lewis \& Short): to feel together. English \emph{consent} comes from it.
\end{cousinweb}

\begin{grammarlens}[title={What you've built}]
One more word for telling, explaining and describing at B1.
\end{grammarlens}

\subsection*{Your turn}
\begin{itemize}
\raggedright
  \item \emph{Say it:} \emph{cōnsentiō, cōnsentīre}
  \item \emph{Say it:} \emph{cōnsentiō, cōnsentīre}, once more
  \item \emph{Say it:} a sentence of your own with \emph{cōnsentiō, cōnsentīre}
  \item \emph{Recall:} say the Latin for to agree, then the Latin for to disagree, then say \emph{cōnsentiō, cōnsentīre} again
\end{itemize}

\begin{tcolorbox}[breakable,colback=teal!4,colframe=teal!35!black,title={Before you move on}]
What does cōnsentiō, cōnsentīre mean? (\textquotedblleft{}To agree together, to be of one mind\textquotedblright{}.) Say it once more.
\end{tcolorbox}
\section[cōnsēnsus, cōnsēnsūs]{cōnsēnsus, cōnsēnsūs — agreement, unanimity}

\label{lesson:LA-C358-consensus}



\begin{quote}
Before the new one: say the Latin for to disagree, then the Latin for to agree together.
\end{quote}

\subsection*{You'll want to know: cōnsēnsus, cōnsēnsūs}
\textbf{cōnsēnsus, cōnsēnsūs} — \textquotedblleft{}agreement, unanimity\textquotedblright{}.

\begin{quote}
Omnium cōnsēnsū eum ducem fēcērunt.
\end{quote}

\begin{quote}
\emph{By common consent they made him leader.}
\end{quote}

\emph{cōnsēnsus, cōnsēnsūs} m., fourth declension. \emph{Omnium cōnsēnsū}, by everyone's agreement, is a stock phrase. English \emph{consensus} is the same word.

\begin{cousinweb}
\emph{cōnsēnsus} is the noun built on \emph{cōnsentiō} \textquotedblleft{}to agree together\textquotedblright{} (EtymDB; Lewis \& Short), and English took it over whole as \emph{consensus}.
\end{cousinweb}

\begin{grammarlens}[title={What you've built}]
One more word for telling, explaining and describing at B1.
\end{grammarlens}

\subsection*{Your turn}
\begin{itemize}
\raggedright
  \item \emph{Say it:} \emph{cōnsēnsus, cōnsēnsūs}
  \item \emph{Say it:} \emph{cōnsēnsus, cōnsēnsūs}, once more
  \item \emph{Say it:} a sentence of your own with \emph{cōnsēnsus, cōnsēnsūs}
  \item \emph{Recall:} say the Latin for to disagree, then the Latin for to agree together, then say \emph{cōnsēnsus, cōnsēnsūs} again
\end{itemize}

\begin{tcolorbox}[breakable,colback=teal!4,colframe=teal!35!black,title={Before you move on}]
What does cōnsēnsus, cōnsēnsūs mean? (\textquotedblleft{}Agreement, unanimity\textquotedblright{}.) Say it once more.
\end{tcolorbox}
\section[contrōversia, contrōversiae]{contrōversia, contrōversiae — a dispute, a controversy}

\label{lesson:LA-C358-controversia}



\begin{quote}
Before the new one: say the Latin for to agree together, then the Latin for agreement.
\end{quote}

\subsection*{You'll want to know: contrōversia, contrōversiae}
\textbf{contrōversia, contrōversiae} — \textquotedblleft{}a dispute, a controversy\textquotedblright{}.

\begin{quote}
Inter vīcīnōs contrōversia dē agrō erat.
\end{quote}

\begin{quote}
\emph{There was a dispute between the neighbours about a field.}
\end{quote}

\emph{contrōversia, contrōversiae} f. A quarrel at law over a boundary or a will, and a set debate in the schools of speaking. English \emph{controversy}.

\begin{cousinweb}
\emph{contrōversia} is built on \emph{contrōversus} \textquotedblleft{}turned the other way\textquotedblright{} (EtymDB; Lewis \& Short), from \emph{contrā} \textquotedblleft{}against\textquotedblright{} and the participle of \emph{vertō} \textquotedblleft{}to turn\textquotedblright{}: two sides turned against each other. English \emph{controversy} comes from it.
\end{cousinweb}

\begin{grammarlens}[title={What you've built}]
One more word for telling, explaining and describing at B1.
\end{grammarlens}

\subsection*{Your turn}
\begin{itemize}
\raggedright
  \item \emph{Say it:} \emph{contrōversia, contrōversiae}
  \item \emph{Say it:} \emph{contrōversia, contrōversiae}, once more
  \item \emph{Say it:} a sentence of your own with \emph{contrōversia, contrōversiae}
  \item \emph{Recall:} say the Latin for to agree together, then the Latin for agreement, then say \emph{contrōversia, contrōversiae} again
\end{itemize}

\begin{tcolorbox}[breakable,colback=teal!4,colframe=teal!35!black,title={Before you move on}]
What does contrōversia, contrōversiae mean? (\textquotedblleft{}A dispute, a controversy\textquotedblright{}.) Say it once more.
\end{tcolorbox}
